\textcolor{red}{\textbf{La primera respuesta generada por la IA produjo un bloque \texttt{flowchart LR} sintácticamente válido, que renderizó sin errores en mermaid.live y respetaba las reglas de la plantilla: cabecera \texttt{flowchart LR}, \texttt{subgraph} como límite del sistema, casos de uso como nodos redondeados, actor con icono y salto de línea, una única asociación actor--caso de uso, relaciones «include» con la dirección correcta, estereotipos con comillas latinas y \texttt{classDef}/\texttt{class} para ocultar el borde del actor. No obstante, tras auditar el contenido frente a los criterios de aceptación de US-1 y al resto del modelo del proyecto, se aplicaron las siguientes correcciones y mejoras:
\begin{itemize}[leftmargin=*,nosep]
\item \textbf{Caso de uso «Calcular precio con IVA» incoherente con el modelo de datos:} La IA planteó el IVA como un cálculo del sistema. Sin embargo, el criterio de aceptación exige \emph{mostrar} el precio con IVA y el catálogo ya almacena ese valor (atributo \texttt{precio\_iva} de la entidad \texttt{PRODUCTO}). Se sustituyó por «Mostrar imagen, nombre y precio con IVA».
\item \textbf{Caso de uso ambiguo que solapa con US-2:} «Mostrar ficha básica de producto» se confunde con la «Ficha de detalle de producto» de US-2 y no recoge explícitamente la imagen ni el nombre exigidos por los criterios. Se reformuló para reflejar exactamente los datos de la tarjeta del catálogo y se mantuvo «Mostrar distintivo de disponibilidad» como caso incluido propio.
\item \textbf{Ausencia de relaciones «extend» (flujos condicionales):} La respuesta solo contenía relaciones «include». Se añadieron dos flujos opcionales: «Mostrar mensaje de catálogo vacío» extiende a «Ver catálogo principal» (coherente con la rama \emph{Catálogo vacío} del diagrama de secuencia de US-1) y «Mostrar aviso de producto agotado» extiende a «Mostrar distintivo de disponibilidad».
\item \textbf{Identificadores y título poco descriptivos:} Los identificadores \texttt{CU\_US1\_1}, \texttt{CU\_US1\_2} y \texttt{CU\_US1\_3} no indican su función y el título del límite del sistema era genérico. Se renombraron con identificadores descriptivos (por ejemplo, \texttt{MostrarDatos}) y el límite pasó a titularse «Atelier Jaén - Catálogo de la tienda», con el caso principal alineado con el título de la historia («Ver catálogo principal»).
\end{itemize}}}

Describid aquí los errores \textbf{reales} de la primera respuesta de la IA para el caso de uso de US-1 y cómo los corrigió el equipo. Puntos que hay que comprobar:
\begin{itemize}[leftmargin=*,nosep]
\item Más de una asociación actor-caso de uso (solo se permite una).
\item Flecha include en sentido contrario (debe ir del caso base al incluido).
\item Flecha extend en sentido contrario (debe ir del opcional al base) o no discontinua.
\item Estereotipos escritos con < > o << >> en lugar de « ».
\item Falta classDef/class del actor o falta el icono fa:fa-user.
\item ID del subgraph repetido como ID de nodo; paréntesis en etiquetas.
\item Casos de uso inventados que no aparecen en los criterios de aceptación, o criterios sin reflejar.
\end{itemize}
